\chapter{Right Triangle Trigonometry}
\label{ch:ch03-right-triangle}

Chapter~1 established that certain ratios survive changes of scale. Chapter~2
established vectors, and showed that every vector generates a right
triangle whose legs are its components. This chapter asks the question that
those two results were building toward: given a right triangle with a fixed
acute angle, what remains unchanged as the triangle is enlarged or reduced?
The answer will turn out to be a small, fixed collection of ratios, and
naming those ratios is all that defining sine, cosine, and tangent amounts
to.

\section{Motivation: The Problem of Direction}

A vector $\mathbf{v} = \langle x,y \rangle$ has a magnitude,
$|\mathbf{v}| = \sqrt{x^2+y^2}$, computed in the previous chapter. Magnitude
alone, however, does not determine a vector: a vector of length $10$ may
point in infinitely many directions, and nothing established so far can
distinguish one from another. Measuring direction will require a new idea,
and that idea begins with the right triangle every vector already carries
inside it.

\section{Similar Right Triangles}

Construct several right triangles sharing the same acute angle $\theta$: for
instance, triangles with sides $3$-$4$-$5$, $6$-$8$-$10$, and $9$-$12$-$15$.
Each is obtained from the $3$-$4$-$5$ triangle by scaling, so by
Angle--Angle similarity all three triangles are similar, and Chapter~1
guarantees that corresponding side ratios agree across all of them:
\[
\frac{3}{5} = \frac{6}{10} = \frac{9}{15}, \qquad
\frac{4}{5} = \frac{8}{10} = \frac{12}{15}.
\]

\begin{theorem}
For any collection of right triangles sharing a common acute angle $\theta$,
the ratio of any two corresponding sides is the same across every triangle
in the collection.
\end{theorem}

\begin{proof}
Two right triangles sharing an acute angle $\theta$ also share their right
angle, hence agree in two corresponding angles and are similar by the
Angle--Angle criterion of Chapter~1. Similar triangles have proportional
corresponding sides, so the ratio of any two corresponding sides is the
common scale factor's cancellation, exactly as shown in the derivation of
Section~1.5: the scale factor appears in both the numerator and denominator
of the ratio and cancels, leaving a value depending only on $\theta$ and not
on the size of the triangle chosen.
\end{proof}

This is the central theorem of the chapter, and everything that follows is
simply a matter of giving names to the ratios it guarantees exist.

\section{Naming the Invariants}

For an acute angle $\theta$ in a right triangle, call the side across from
$\theta$ the \emph{opposite} side, the side adjacent to $\theta$ that is not
the hypotenuse the \emph{adjacent} side, and the side across from the right
angle the \emph{hypotenuse}.

\begin{definition}
For an acute angle $\theta$ of a right triangle,
\[
\sin\theta = \frac{\text{opposite}}{\text{hypotenuse}}, \qquad
\cos\theta = \frac{\text{adjacent}}{\text{hypotenuse}}, \qquad
\tan\theta = \frac{\text{opposite}}{\text{adjacent}}.
\]
\end{definition}

\begin{historicalnote}{Where the Word ``Sine'' Comes From}
The word \emph{sine} has an unusually crooked history for a term this
central to mathematics \cite{boyer1991}. Indian astronomers of the first
millennium worked with a quantity called \emph{jy\={a}-ardha}, literally
``half-chord,'' shortened in later Sanskrit texts to \emph{jy\={a}} or
\emph{jiv\={a}}. When these texts were translated into Arabic, the word
was transliterated phonetically as \emph{jiba}, a term with no meaning of
its own in Arabic. Arabic is ordinarily written without vowels, so later
readers encountering the consonants of \emph{jiba} reinterpreted them as
the ordinary Arabic word \emph{jaib}, meaning a bay, fold, or pocket in a
garment. When the twelfth-century translator Gerard of Cremona rendered
this word into Latin, he translated the meaning of \emph{jaib} rather
than transliterating the original Sanskrit, giving the Latin word
\emph{sinus}, meaning a bay, curve, or fold of cloth, from which English
inherited \emph{sine}. The name of one of the most important functions in
mathematics therefore descends from a fold in a piece of clothing, by way
of a translator's honest misreading of an unfamiliar word.
Cosine and cotangent follow more transparently: the prefix ``co-''
abbreviates \emph{complementi}, since the cosine of $\theta$ is the sine
of its complementary angle, $90^\circ - \theta$.
\end{historicalnote}

It is worth pausing on what has actually happened here. These three
equations are not arbitrary formulas handed down by convention; they are
names attached to quantities the theorem of the previous section has already
proved to be invariant under scaling. Nothing about $\sin\theta$,
$\cos\theta$, or $\tan\theta$ depends on which triangle with acute angle
$\theta$ happens to be drawn, precisely because the ratio defining each one
was shown, in general, to be scale-independent before it was ever given a
name.

\begin{sidebar}{Why Do We Name Invariants?}
A quantity earns a name in mathematics and science largely because it
survives changes that are, for the purpose at hand, irrelevant. Temperature
is useful because it survives the shape of its container; density is useful
because it survives the size of the sample measured; a trigonometric ratio
is useful because it survives the scale of the triangle it is computed from.
In each case a name is attached not to a number but to an invariant, a
quantity that stays fixed while everything else around it changes, and it is
this stability that makes the quantity worth comparing across situations in
the first place. Mathematics repeatedly performs this same maneuver ---
identify what a family of transformations leaves unchanged, then name it ---
and sine, cosine, and tangent are among the earliest and most durable
examples a student will encounter.
\end{sidebar}

\section{The Remaining Three Functions}

Three further ratios are traditionally defined as reciprocals of the ones
above:
\[
\csc\theta = \frac{1}{\sin\theta}, \qquad
\sec\theta = \frac{1}{\cos\theta}, \qquad
\cot\theta = \frac{1}{\tan\theta}.
\]
These reciprocal functions --- cosecant, secant, and cotangent --- appear
throughout the historical literature and remain useful in certain identities
and calculus formulas encountered later in this book, but the overwhelming
majority of applications in this text and elsewhere rely on sine, cosine,
and tangent alone.

\section{Finding Unknown Sides}

\begin{algorithmproc}{How to Find an Unknown Side Using a Trigonometric Ratio}
Given a right triangle with one known acute angle and one known side, first
identify which of the two known quantities the unknown side is related to
by determining whether the unknown side is opposite, adjacent, or the
hypotenuse relative to the known angle. Next, select whichever of
$\sin\theta$, $\cos\theta$, or $\tan\theta$ relates the known side, the
unknown side, and the known angle. Write the resulting equation and solve it
for the unknown side.
\end{algorithmproc}

\begin{example}
A ten-meter ladder leans against a wall, making a $65^\circ$ angle with the
ground. How high up the wall does the ladder reach? The height $h$ is
opposite the given angle, and the ladder itself is the hypotenuse, so
\[
\sin(65^\circ) = \frac{h}{10}, \qquad h = 10\sin(65^\circ) \approx 9.06
\text{ meters.}
\]
\end{example}

\begin{practicenow}
A twelve-meter ladder leans against a wall, making a $72^\circ$ angle with
the ground. Find the height the ladder reaches up the wall.
\end{practicenow}

\section{Finding Unknown Angles}

The reverse problem also arises naturally: given the ratio of two sides, how
is the corresponding angle recovered? If $\tan\theta = 7/4$, for example,
some process must undo the tangent function and return $\theta$ itself. Such
a process exists and is used freely at this stage with a calculator, but it
will not be formally defined until Chapter~7 introduces the inverse
trigonometric functions with the care that definition deserves. For now it
is enough to know that the question is meaningful, and that a numerical
answer can be obtained.

\section{Special Right Triangles}

Two right triangles recur so often that their ratios are worth deriving
explicitly, once and for all, rather than looked up in a table.

\begin{example}[The $45^\circ$--$45^\circ$--$90^\circ$ Triangle]
Consider an isosceles right triangle with both legs of length $1$. By the
Pythagorean theorem, its hypotenuse has length $\sqrt{1^2+1^2} = \sqrt{2}$,
giving side lengths $1$, $1$, $\sqrt{2}$. Since the two legs are equal, the
two acute angles are also equal, and since they sum to $90^\circ$ each must
equal $45^\circ$. Applying the definitions of the previous section,
\[
\sin 45^\circ = \cos 45^\circ = \frac{1}{\sqrt{2}} = \frac{\sqrt2}{2}, \qquad
\tan 45^\circ = \frac{1}{1} = 1.
\]
\end{example}

\begin{example}[The $30^\circ$--$60^\circ$--$90^\circ$ Triangle]
Consider an equilateral triangle with side length $2$, so that all three
angles equal $60^\circ$. Dropping an altitude from one vertex to the
midpoint of the opposite side bisects both that side and the angle at the
vertex, producing two congruent right triangles, each with hypotenuse $2$,
short leg $1$ (half of the original side), and a $30^\circ$--$60^\circ$--$90^\circ$
angle configuration. By the Pythagorean theorem the remaining leg has length
$\sqrt{2^2 - 1^2} = \sqrt{3}$, giving side lengths $1$, $\sqrt3$, $2$.
Applying the definitions of the previous section to the $30^\circ$ angle,
whose opposite side is the short leg of length $1$ and whose adjacent side
is the leg of length $\sqrt3$,
\[
\sin 30^\circ = \frac{1}{2}, \qquad
\cos 30^\circ = \frac{\sqrt3}{2}, \qquad
\tan 30^\circ = \frac{1}{\sqrt3} = \frac{\sqrt3}{3}.
\]
The complementary angle $60^\circ$ exchanges the roles of opposite and
adjacent, giving $\sin 60^\circ = \sqrt3/2$, $\cos 60^\circ = 1/2$, and
$\tan 60^\circ = \sqrt3$.
\end{example}

\section{Applications}

Right-triangle trigonometry is the working tool of surveying, where an
inaccessible height or distance is recovered from an angle of elevation or
depression together with a measured baseline; of navigation, where a
vessel's position is found from bearing and distance; and of architecture
and engineering, where the ratios of this chapter govern everything from
roof pitch to the forces along a diagonal brace. The mathematics developed
above is what makes each of these applications possible, and the
applications themselves are best treated as exercises in identifying which
side and angle are known and which ratio connects them.

\section{Looking Ahead}

This chapter has defined the trigonometric functions using right triangles,
and that definition carries a built-in restriction worth stating plainly.
An acute angle inside a right triangle can never reach or exceed $90^\circ$,
and it can never be negative, so as it stands sine, cosine, and tangent are
only defined for a narrow range of angles. If trigonometry is to become a
general theory of direction, capable of describing a vector pointing in any
direction whatsoever, this restriction must eventually be removed. The next
chapter develops a framework capable of describing any angle at all, by
returning to the unit circle introduced informally in Chapter~0 and asking
what a trigonometric ratio can mean once the right triangle that first
defined it is allowed to disappear.

The pattern established in this chapter, transformation to invariant to
name, is one that recurs throughout the book: a family of transformations
(here, scaling) is shown to preserve some quantity (here, a ratio of sides),
and that preserved quantity is given a name (here, sine, cosine, and
tangent) precisely because its invariance is what makes it worth naming.
The same pattern will reappear with rotations in Chapter~8, with complex
numbers in Chapter~16, and with Fourier modes in Chapter~19.

\section{Exercises}
\begin{exercise}
A right triangle has an acute angle of $38^\circ$ and a hypotenuse of $14$
centimeters. Find the lengths of the two legs.
\end{exercise}

\begin{exercise}
Derive the values of $\sin 45^\circ$, $\cos 45^\circ$, and $\tan 45^\circ$
directly from a right triangle with legs of length $2$ instead of $1$, and
confirm that the result agrees with the derivation given in the text.
\end{exercise}

\begin{exercise}
A surveyor stands $40$ meters from the base of a tower and measures the
angle of elevation to the top of the tower as $52^\circ$. Find the height
of the tower.
\end{exercise}

\begin{exercise}
Explain, using the theorem of Section~3.2, why $\sin\theta$ could not
possibly depend on which right triangle with acute angle $\theta$ is used to
compute it.
\end{exercise}
